\section{Results}
\label{sec:results}

\subsection{Executability and bounded restoration}
\label{sec:results-executability}

Primary execution evaluated 13 of the 39 frozen cases (33.3\%), while 26
cases (66.7\%) were non-evaluable at the primary layer. These non-evaluable
cases were retained in the study frame rather than replaced by additional
repositories.

\begin{table}[t]
\centering
\caption{Executability before and after bounded restoration. Non-evaluable cases remain part of the frozen 39-case study frame.}
\label{tab:study-summary}
\vspace{0.35em}
\small
\renewcommand{\arraystretch}{1.28}
\setlength{\tabcolsep}{10pt}
\begin{tabular}{lrrr}
\toprule
\addlinespace[2pt]
\textbf{Layer} & \textbf{Selected} & \textbf{Evaluated, n (\%)} & \textbf{Non-evaluable, n (\%)} \\
\addlinespace[2pt]
\midrule
\addlinespace[3pt]
Primary execution & 39 & 13 (33.3\%) & 26 (66.7\%) \\
\addlinespace[3pt]
\textbf{Combined primary + D027} & 39 & \textbf{24 (61.5\%)} & \textbf{15 (38.5\%)} \\
\addlinespace[3pt]
\bottomrule
\end{tabular}
\vspace{0.25em}
\end{table}


The bounded D027 restoration layer increased the number of evaluable cases
from 13 to 24, corresponding to 61.5\% of the complete frozen frame. Fifteen
of 39 cases (38.5\%) remained non-evaluable after combining primary and
restoration evidence.

Within the 24-case D027 cohort, a report or assessment was present for 19
cases. Substantive restoration was attempted for 18 cases. Eleven cases were
successfully restored and evaluated, while seven were not restored after a
substantive attempt. One additional case had a D027 report but no substantive
restoration attempt, and five cohort cases had no D027 report.

These counts characterize executability under the study protocol rather than
mutation detection. In particular, setup and workflow failures were not
assigned KILLED or SURVIVED outcomes.


\subsection{RQ1: Detection of reproducibility-relevant mutations}
\label{sec:results-rq1}

Among the 24 combined evaluated cases, one mutation was confirmed
semantically equivalent and was excluded from the meaningful detection
denominator. The resulting denominator contained 23 confirmed
non-equivalent mutations.

Of these 23 mutations, 2 were \textsc{Killed} and 21 were
\textsc{Survived}. The observed detection proportion was therefore

\[
\frac{2}{23} = 8.7\%.
\]

Thus, in this sample, the selected repository validation workflows detected
2 of 23 confirmed non-equivalent reproducibility-relevant mutations. The
remaining 21 confirmed non-equivalent mutations completed without detection
by their selected frozen workflows.

The B02 operator-specific results were uneven. All five meaningful
\texttt{random-seed} evaluations survived. For \texttt{dependency-pin},
four cases were evaluated: two were KILLED, one was SURVIVED, and one was
confirmed equivalent. Excluding the equivalent case yielded two detections
among three confirmed non-equivalent \texttt{dependency-pin} evaluations.
All four meaningful \texttt{data-split} cases and all three
\texttt{cv-fold-count} cases survived.

\begin{table}[t]
\centering
\caption{Combined B02 results by mutation operator. Detection is reported over confirmed non-equivalent evaluated mutations; confirmed-equivalent evaluations are excluded.}
\label{tab:operator-results}
\vspace{0.35em}
\small
\renewcommand{\arraystretch}{1.24}
\setlength{\tabcolsep}{6pt}
\begin{tabular}{lrrrrrr}
\toprule
\addlinespace[2pt]
\textbf{Operator} & \textbf{Eval./sel.} & \textbf{Non-eval.} & \textbf{Killed} & \textbf{Surv.} & \textbf{Eq.} & \textbf{Detection} \\
\addlinespace[2pt]
\midrule
\addlinespace[3pt]
\texttt{random-seed} & 5/10 & 5 & 0 & 5 & 0 & 0/5 (0.0\%) \\
\addlinespace[3pt]
\texttt{dependency-pin} & 4/10 & 6 & \textbf{2} & 1 & 1 & \textbf{2/3 (66.7\%)} \\
\addlinespace[3pt]
\texttt{data-split} & 4/6 & 2 & 0 & 4 & 0 & 0/4 (0.0\%) \\
\addlinespace[3pt]
\texttt{cv-fold-count} & 3/3 & 0 & 0 & 3 & 0 & 0/3 (0.0\%) \\
\addlinespace[3pt]
\bottomrule
\end{tabular}
\vspace{0.25em}
\end{table}


These operator-level counts are descriptive. The operator-specific
denominators are small and unequal, and the corpus construction was not
designed to estimate population-level differences in mutation detection
between operators.


\subsection{RQ2: Detection by workflow and oracle type}
\label{sec:results-rq2}

We first stratified the 23 confirmed non-equivalent evaluated mutations by the
frozen validation-workflow category.

Among six \texttt{upstream-test} cases, none of the mutations was detected.
The single CI case also survived. One of two
\texttt{documented-validation} cases was detected, as was one of eleven
\texttt{documented-example} cases. None of the three
\texttt{documented-experiment} cases was detected.

\begin{table}[t]
\centering
\caption{Detection of confirmed non-equivalent mutations by frozen validation-workflow type. Workflow type is categorical; no ordinal ranking is assumed.}
\label{tab:rq2-workflow}
\begin{tabular}{lrrrr}
\toprule
Workflow type & N & Killed & Survived & Detection \\
\midrule
upstream-test & 6 & 0 & 6 & 0.0\% \\
CI & 1 & 0 & 1 & 0.0\% \\
documented-validation & 2 & 1 & 1 & 50.0\% \\
documented-experiment & 3 & 0 & 3 & 0.0\% \\
documented-example & 11 & 1 & 10 & 9.1\% \\
\bottomrule
\end{tabular}
\end{table}


The primary oracle analysis was restricted to B02 because oracle kind was
prospectively recorded only under protocol version 2. Fifteen B02 cases
entered the confirmed non-equivalent meaningful denominator. Of these, two
used an \texttt{assertion} oracle and thirteen used a
\texttt{completion-only} oracle.

Neither of the two \texttt{assertion} cases was detected. Two of the thirteen
\texttt{completion-only} cases were detected, corresponding to a descriptive
detection proportion of 15.4\%.

\begin{table}[t]
\centering
\caption{Detection by prospectively recorded validation-oracle kind in B02. B01 is excluded because schema version 1 did not prospectively record oracle kind.}
\label{tab:rq2-oracle}
\begin{tabular}{lrrrr}
\toprule
Oracle kind & N & Killed & Survived & Detection \\
\midrule
assertion & 2 & 0 & 2 & 0.0\% \\
completion-only & 13 & 2 & 11 & 15.4\% \\
\bottomrule
\end{tabular}
\end{table}


This pattern does not establish that completion-only workflows provide
stronger mutation detection than assertion-based workflows. A diagnostic
cross-classification by operator showed that both detected B02 mutations
occurred in the \texttt{dependency-pin} $\times$
\texttt{completion-only} cell: two of the three confirmed non-equivalent
cases in that cell were KILLED. No KILLED mutation occurred in any other
observed B02 operator--oracle cell.

The oracle categories also differed in evaluability. Of seven prospectively
classified \texttt{assertion} cases, three were combined-evaluable and two
entered the confirmed non-equivalent meaningful denominator. Of 22
\texttt{completion-only} cases, thirteen were combined-evaluable and all
thirteen entered the meaningful denominator.

Together, the sparse cells, operator composition, and differential
evaluability prevent attributing the observed detection pattern to oracle
strength. RQ2 therefore provides descriptive evidence of variation across
workflow and oracle categories, but no evidence of a monotonic relationship
between oracle explicitness and mutation detection.
